\titleformat{\chapter}[block]{\Huge\bfseries}{
    \sffamily\huge\centering{Appendice \thechapter
}\\}{1ex}{\sffamily\centering\huge\bfseries}

\newpage\pagestyle{appendices}
\begin{appendices}

\chapter{Codice Sorgente}\label{app:codice_sorgente}

In questa appendice viene riportato il codice sorgente in linguaggio Python relativo allo step di estrazione dei documenti PDF, con particolare focus sull'implementazione del \textit{Custom Triage Processor}. Il sistema è stato progettato adottando lo \textit{Strategy Pattern}, che permette di definire una famiglia di algoritmi, incapsularli e renderli intercambiabili in base alla natura del documento analizzato.

\section{Interfaccia Base}
Il file \texttt{base.py} definisce l'interfaccia astratta che tutte le strategie di elaborazione dei PDF devono implementare.

\begin{lstlisting}[language=Python, breaklines=true, escapeinside={}, caption={Interfaccia base per le strategie}, label={lst:base_strategy}]
class PDFProcessingStrategy(ABC):
    @abstractmethod
    def process(self, pdf_file: Path, reader: PdfReader, temp_export_path: Path) -> str:
        pass
\end{lstlisting}

\section{Implementazione delle Strategie}

\subsection{Strategia per Tagged PDF}
Il file \texttt{tagged\_strategy.py} implementa la strategia per i documenti nativamente strutturati.

\begin{lstlisting}[language=Python, breaklines=true, escapeinside={}, caption={Strategia per PDF strutturati}, label={lst:tagged_strategy}]
class TaggedPDFStrategy(PDFProcessingStrategy):
    def process(self, pdf_file: Path, reader: PdfReader, temp_export_path: Path) -> str:
        final_markdown = ""
        convert_args = {
            "input_path": [str(pdf_file)],
            "output_dir": str(temp_export_path),
            "format": ["markdown"],
            "use_struct_tree": True,
            "hybrid_mode": "auto"
        }
        
        try:
            opendataloader_pdf.convert(**convert_args)
            for temp_md_file in temp_export_path.glob("*.md"):
                with open(temp_md_file, "r", encoding="utf-8") as f:
                    final_markdown += f.read()
                temp_md_file.unlink()
        except Exception as e:
            logger.error(f"Errore durante l'estrazione di {pdf_file.name}: {e}")
        
        return final_markdown
\end{lstlisting}

\subsection{Custom Triage Processor}
Il file \texttt{triage\_strategy.py} contiene il core della logica di smistamento intelligente delle singole pagine.

\begin{lstlisting}[language=Python, breaklines=true, escapeinside={}, caption={Implementazione del Custom Triage Processor}, label={lst:triage_strategy}]
class TriagePDFStrategy(PDFProcessingStrategy):
    STRATEGY_MAP = {
        PageType.SIMPLE: {"hybrid_mode": "auto"},
        PageType.COMPLEX: {"hybrid_mode": "full", "hybrid": "docling-fast", "hybrid_url": URL_DOCLING_NO_OCR},
        PageType.SCAN: {"hybrid_mode": "full", "hybrid": "docling-fast", "hybrid_url": URL_DOCLING_OCR}
    }

    def process(self, pdf_file: Path, reader: PdfReader, temp_export_path: Path) -> str:
        final_markdown = ""

        for i, page in enumerate(reader.pages):
            page_strategy = self._analyze_page_type(page)
            
            temp_pdf_path = temp_export_path / f"temp_p{i}.pdf"
            with open(temp_pdf_path, "wb") as f:
                writer = PdfWriter()
                writer.add_page(page)
                writer.write(f)

            convert_args = {
                "input_path": [str(temp_pdf_path)],
                "output_dir": str(temp_export_path),
                "format": ["markdown"],
                "use_struct_tree": False
            }

            convert_args.update(self.STRATEGY_MAP.get(page_strategy, self.STRATEGY_MAP[PageType.SIMPLE]))

            try:
                opendataloader_pdf.convert(**convert_args)
                for temp_md_file in temp_export_path.glob("*.md"):
                    with open(temp_md_file, "r", encoding="utf-8") as f:
                        final_markdown += f.read() + "\n\n---\n\n"
                    temp_md_file.unlink()
            except Exception as e:
                logger.error(f"Errore critico pagina {i+1} di {pdf_file.name}: {e}")
            finally:
                if temp_pdf_path.exists():
                    temp_pdf_path.unlink()
                
        return final_markdown

    def _analyze_page_type(self, page) -> PageType:
        text = page.extract_text() or ""
        has_images = len(page.images) > 0 if hasattr(page, 'images') else False

        if len(text.strip()) < 50 and has_images:
            return PageType.SCAN
        if has_images or len(text.strip()) > 2500:
            return PageType.COMPLEX

        return PageType.SIMPLE
\end{lstlisting}

\section{Logica di Selezione della Strategia}
Il file \texttt{utils.py} funge da \textit{Factory} per il pattern architetturale scelto: analizza i metadati del documento per verificare la presenza di un albero strutturale e istanzia dinamicamente la strategia di elaborazione più adeguata.

\begin{lstlisting}[language=Python, breaklines=true, escapeinside={}, caption={Factory per la selezione della strategia}, label={lst:utils}]
def get_pdf_strategy(pdf_reader) -> PDFProcessingStrategy:
    if is_pdf_tagged(pdf_reader):
        return TaggedPDFStrategy()
    return TriagePDFStrategy()


    
def is_pdf_tagged(reader: PdfReader) -> bool:
    try:
        return "/StructTreeRoot" in reader.trailer["/Root"]
    except Exception:
        return False
\end{lstlisting}

\section{Esecuzione Principale}
Infine, lo script \texttt{main.py} coordina il processo, determinando quale strategia applicare.

\begin{lstlisting}[language=Python, breaklines=true, escapeinside={}, caption={Script principale di esecuzione}, label={lst:main}]
def main():
    setup_logging(LOG_DIR)
    logger = logging.getLogger(__name__)
    
    for path in [INPUT_PDF_DIR, BASE_MD_PATH, TEMP_EXPORT_PATH]:
        path.mkdir(parents=True, exist_ok=True)

    for pdf_file in INPUT_PDF_DIR.iterdir():                    
        try:
            pdfReader = PdfReader(pdf_file)
            pdf_extraction_strategy = get_pdf_strategy(pdfReader)
            
            markdown = pdf_extraction_strategy.process(pdf_file, pdfReader, TEMP_EXPORT_PATH)

            with open(BASE_MD_PATH / f"{pdf_file.stem}.md", "w", encoding="utf-8") as f:
                f.write(markdown)
        except Exception as e:
            logger.error(f"Errore su {pdf_file.name}: {e}")

    if TEMP_EXPORT_PATH.exists():
        shutil.rmtree(TEMP_EXPORT_PATH)

if __name__ == "__main__":
    main()
\end{lstlisting}

\end{appendices}